
Der Strategie-Fit beschreibt die Passung zwischen den langfristigen Zielen
eines KMU und dem Zweck, Umfang und erwarteten Nutzen einer ERP-Einführung.

Vor der Einführung muss ein KMU zunächst klären, welchen strategischen Zweck
das ERP-System erfüllen soll \parencite{Leyh2015}. Dazu gehört die Formulierung einer Business
Vision und eines Business Case, in dem der erwartete Nutzen des Systems
konkret beschrieben wird \parencite{Leyh2015}. Strategische Ziele können etwa die Stärkung der
Wettbewerbsfähigkeit, die Verbesserung der Unternehmenssteuerung, die
Senkung von Kosten, die Erhöhung der Skalierbarkeit oder die Vorbereitung
weiterer Digitalisierungsschritte sein \parencite{Leyh2015}. Entscheidend ist, dass diese Ziele
nicht nur allgemein genannt, sondern vor der Systemauswahl in Anforderungen
und Prioritäten übersetzt und konkretisiert werden \parencite{Jha2008,Leyh2015}.

Strategie-Fit bedeutet damit auch, dass ein ERP-System nicht als isoliertes
IT-Projekt verstanden wird. Die Einführung verändert die Art und Weise, wie das Unternehmen Prozesse steuert,
Informationen nutzt und Entscheidungen trifft \parencite{Leyh2017}.

Eine zentrale Rolle übernimmt dabei die Geschäftsführung. Sie muss die
strategische Bedeutung des ERP-Projekts frühzeitig bestimmen, Ziele
priorisieren und die Einführung als unternehmensweites Vorhaben legitimieren \parencite{Leyh2015}.
Entscheidungen über Projektumfang, Systemauswahl,
Standardisierung, Anpassung und erwarteten Nutzen dürfen nicht allein auf die
IT-Ebene verlagert werden \parencite{Leyh2015}.

Strategie-Fit zeigt sich außerdem in der Auswahl des passenden Systems. Ein
ERP-System muss nicht nur aktuelle Anforderungen erfüllen, sondern auch zur
zukünftigen Entwicklung des KMU passen. Dazu gehören Skalierbarkeit,
Anpassungsfähigkeit und die Fähigkeit, künftige Digitalisierungsinitiativen
zu unterstützen. Für KMU ist dies besonders relevant, weil ihre
Wettbewerbsfähigkeit häufig auf Flexibilität und schneller Anpassung an
Marktveränderungen beruht. Ein ERP-System darf diese Agilität nicht durch
zu starre Standardprozesse einschränken, sondern sollte sie im Rahmen einer
integrierten Unternehmenssteuerung unterstützen \parencite{Leyh2015,Leyh2017}.


Ein hoher Strategie-Fit liegt somit vor, wenn das KMU vor der
ERP-Einführung klar definiert hat, welchen langfristigen, strategischen Beitrag das System
leisten soll. Dazu gehören eine dokumentierte Zielsetzung, ein realistischer
Business Case, eine Verbindung zur Digitalisierungsstrategie, eine aktive
strategische Steuerung durch die Geschäftsführung und eine Systemauswahl,
die zur langfristigen Entwicklung des Unternehmens passt. 
